% ============================================================================
% Fichier  : annexes/Ann_B_outils.tex
% Titre    : Référentiel des Outils
% Version  : 1.0 -- novembre 2024
% ============================================================================

\chapter{Référentiel des Outils}
\label{ann:outils}

Les outils cités dans le cours sont rangés par usage. La plupart sont libres
ou gratuits. L'annexe n'est pas une liste d'achats~: un outil ne vaut que par
la méthode de celui qui l'emploie, sa version notée et sa validation sur une
pièce de référence (Chapitre~\ref{chap:normes-droit}, ISO/IEC~27041).

\begin{boxalert}
\textbf{Règle de l'autorisation écrite.} Aucun outil d'acquisition,
d'analyse ou de balayage ne s'emploie sur un système sans l'autorisation
écrite de celui qui en a le droit, et, le cas échéant, de l'hébergeur ou du
prestataire. Sans autorisation, l'usage peut constituer un accès non
autorisé au sens de l'article~65, alinéa~2, de la loi n°~2010/012.
\end{boxalert}

\section*{Acquisition}

\begin{tableaucours}{L{3.6cm}L{5.0cm}L{3.6cm}}{%
    \tableauHeader{Outil} &
    \tableauHeader{Usage} &
    \tableauHeader{Chapitre}
}
Bloqueur d'écriture matériel & Acquisition à froid sans écriture sur la
source & \ref{chap:acquisition} \\
FTK Imager & Images de disques (brut, E01), captures mémoire, calcul
d'empreintes & \ref{chap:acquisition} \\
\texttt{dd}, \texttt{dc3dd}, \texttt{ewfacquire} & Images brutes ou E01 sous
Linux & \ref{chap:acquisition} \\
WinPmem, DumpIt & Capture mémoire sous Windows & \ref{chap:acquisition} \\
LiME, AVML & Capture mémoire sous Linux & \ref{chap:acquisition} \\
\texttt{sha256sum} & Calcul et vérification d'empreintes &
\ref{chap:preuve} \\
\end{tableaucours}

\section*{Analyse}

\begin{tableaucours}{L{3.6cm}L{5.0cm}L{3.6cm}}{%
    \tableauHeader{Outil} &
    \tableauHeader{Usage} &
    \tableauHeader{Chapitre}
}
The Sleuth Kit, Autopsy & Systèmes de fichiers, fichiers supprimés,
artefacts & \ref{chap:disque} \\
Volatility~3 & Analyse de captures mémoire & \ref{chap:memoire} \\
Wireshark, tshark & Analyse de captures réseau & \ref{chap:reseau} \\
Zeek & Journaux par protocole à partir d'une capture & \ref{chap:reseau} \\
log2timeline, Plaso & Super-chronologies & \ref{chap:chronologie} \\
Visionneuses d'événements Windows & Lecture des journaux EVTX &
\ref{chap:disque} \\
\end{tableaucours}

\section*{Vulnérabilités}

\begin{tableaucours}{L{3.6cm}L{5.0cm}L{3.6cm}}{%
    \tableauHeader{Outil ou ressource} &
    \tableauHeader{Usage} &
    \tableauHeader{Chapitre}
}
Scanners de vulnérabilités (libres ou commerciaux) & Balayage authentifié ou
non & \ref{chap:identifier} \\
Calculateur CVSS du FIRST & Score à partir d'un vecteur &
\ref{chap:vulnerabilite} \\
EPSS (FIRST) & Probabilité d'exploitation, interface de consultation &
\ref{chap:vulnerabilite} \\
Catalogue KEV (CISA) & Liste des failles exploitées &
\ref{chap:vulnerabilite}, \ref{chap:gerer} \\
NVD, CVE, CWE & Descriptions et références des failles &
\ref{chap:vulnerabilite} \\
CIS Benchmarks & Références de configuration & \ref{chap:identifier} \\
\end{tableaucours}

\section*{Dans NDAP SUKULU}

Les dix laboratoires du cours n'exigent aucune installation~: leurs pièces
(images réduites, captures, journaux, tableurs de balayage) se lisent avec
les outils du bureau d'analyse de la plateforme ou, pour la version en
ligne, avec un tableur et un éditeur de texte. Chaque laboratoire fournit
son manifeste d'empreintes SHA-256, à vérifier avant tout examen.
